\documentclass[10pt,a4paper]{amsart}
\usepackage[margin=19mm]{geometry}
\usepackage{amsmath,amssymb,mathtools}
\usepackage[scr=rsfs,cal=euler]{mathalfa}
\usepackage{xfrac}
\usepackage[unicode,hidelinks,bookmarks=false]{hyperref}
\newcommand{\ie}{i.e.\ }
\newcommand{\cf}{cf.\ }
\newcommand{\resp}{respectively\ }
\newcommand{\Subsection}{\subsection}
\providecommand{\nobreakdash}{\nobreak\hbox{-}}
\setlength{\emergencystretch}{2em}
\multlinegap0pt
\usepackage{amssymb}
\usepackage{tikz}
\newtheorem{thm}{Theorem}
\title{Local constancy of reduction type and related invariants for curves in $p$-adic families}
\author{Jakab Schrettner}
\date{Source: arXiv:2508.12329v2 (CC BY 4.0)}
\begin{document}
\maketitle

\noindent\emph{Source and licence.} Local constancy of reduction type and related invariants for curves in $p$-adic families. Version arXiv:2508.12329v2 (CC BY 4.0), distributed by arXiv under the Creative Commons Attribution 4.0 International licence, http://creativecommons.org/licenses/by/4.0/, as recorded in the arXiv licence metadata for this exact version and shown on the abstract page https://arxiv.org/abs/2508.12329v2. Full licence notice: \texttt{licenses/arxiv-2508.12329-CC-BY-4.0.txt}.\par\smallskip

\noindent\textit{Editorial scope.} Counted unit: the article's thm Hensel (source0.tex lines 900--907). The article's complete original proof of it is delivered with it (source0.tex lines 908--945, one \texttt{proof} environment of 38 lines). No mathematics is written, repaired or re-derived: the statement and the proof are the article's own lines, in the article's own notation, and no retained line is substituted or re-flowed.  No further source-local context is needed: every cross-reference of every retained line resolves inside the two delivered spans.  Nothing was omitted from any retained span, so the package carries no omission record.  No retained line contains a citation, so the wrapper carries no bibliography and no bibliography entry.  No retained line contains a figure environment, an image-inclusion command or drawing code, so no image is delivered and the package declares no asset dependency.  Editorial formatting only, in addition: an AMS \texttt{amsart} preamble that restates the article's own declarations and macros used by the retained lines, namely the environments \texttt{thm} on separate counters, together with the standard packages those lines need.  The retained environments print in this document as Theorem 1 in order of appearance, whereas the article prints the same items with the numbering of its own full document.  The complete native source of the article as distributed in the arXiv e-print for this exact version is bundled unchanged as \texttt{sources/183394/source0.tex}.
	\begin{thm}[Hensel's lemma]\label{Hensel}
		Let $A$ be a ring, complete with respect to an ideal $I$. Suppose that $f_1, \ldots, f_n\in A[x_1, \ldots, x_n]$ are $n$ polynomials in $n$ variables, and suppose we have $a_1, \ldots, a_n\in A$ such that 
		\[f_i(a_1, \ldots, a_n)\in I^d \text{ for each }i\]
		for some integer $d\ge 1$, and that the $n\times n$ matrix
		\[J = \left(\frac{\partial f_i}{\partial x_j}(a_1, \ldots, a_n)\right)_{i,j = 1}^n\]
		is invertible over $A$. Then there exist $b_1, \ldots, b_n\in A$ such that $b_i\equiv a_i \mod I^d$ and 
		\[f_i(b_1, \ldots, b_n) = 0\text{ for each }i.\]
	\end{thm}

	\begin{proof}
		Define the function $F$ on $A^n$ by
		\[F\left(\left[\begin{array}{c}
			t_1\\
			\vdots\\
			t_n
		\end{array}\right]\right) = \left[\begin{array}{c}
			f_1(t_1, \ldots, t_n)\\
			\vdots\\
			f_n(t_1, \ldots, t_n)
		\end{array}\right].\]
		We change variables as follows: let $g_1, \ldots, g_n\in A[x_1, \ldots, x_n]$ be polynomials such that for $X = (x_1, \ldots, x_n)$ and $\mathbf{a} = (a_1, \ldots, a_n)$ we have
		\[F(J^{-1}(X + \mathbf{a})) = \left[\begin{array}{c}
			g_1(X)\\
			\vdots\\
			g_n(X)
		\end{array}\right] = G(X).\]
		Then by assumption $g_i(0, \ldots, 0) \in I^d$ for all $i$, and the matrix $(\frac{\partial g_i}{\partial x_j}(0, \ldots, 0))$ is the $n\times n$ identity matrix. If we can find $\mathbf{b}' = (b_1', \ldots, b_n')$ such that $g_i(\mathbf{b}') = 0$ for all $i$, and $\mathbf{b}'\in I^dA^n$, then $\mathbf{b} = J\mathbf{b}' + \mathbf{a}$ satisfies the conditions of the theorem. 
		
		Hence we may assume without loss of generality that $\mathbf{a} = (0, \ldots, 0)$ and that $J$ is the identity matrix.
		
		We define the sequence $\mathbf{a}_m$ of elements of $A^n$ by the recurrence
		\[\mathbf{a}_0 = \mathbf{a}\qquad\text{and}\qquad \mathbf{a}_{m+1} = \mathbf{a}_m - F(\mathbf{a}_m).\]
		We claim that $\mathbf{a}_i$ converges to $\mathbf{b} = (b_1, \ldots, b_n)$ which satisfies the requirements of the theorem.
		
		It is clear that for all $m$, $\mathbf{a}_m \in I^dA^n$. We can show by induction that $\mathbf{a}_m \equiv \mathbf{a}_{m+1}\mod I^{d+m}$ for all $m\ge 0$. For $m=0$, this is true by assumption (as $F(\mathbf{a})\in I^dA^n$). Now suppose $m\ge 1$ and the claim holds for all integers up to $m$. We have
		\begin{align*}
			\mathbf{a}_{m+1}-\mathbf{a}_m &= (\mathbf{a}_m - F(\mathbf{a}_m)) - (\mathbf{a}_{m-1} - F(\mathbf{a}_{m-1})) =\\
			&= (\mathbf{a}_m - \mathbf{a}_{m-1}) - (F(\mathbf{a}_m) - F(\mathbf{a}_{m-1})).
		\end{align*}	
		We will use the Taylor expansion
		\[F(X) = F(0) + X + \sum_{i,j=1}^n x_ix_jR_{ij},\]
		which holds for some $R_{ij}\in A^n$ (remember $J$ is the identity). Plugging in $\mathbf{a}_m$ and $\mathbf{a}_{m-1}$, and using $\mathbf{a}_{m}\equiv \mathbf{a}_{m-1} \mod I^{d+m-1}$ we obtain that
		\[F(\mathbf{a}_m)-F(\mathbf{a}_{m-1}) \equiv \mathbf{a}_m - \mathbf{a}_{m-1} \mod I^{2(d+m-1)}\]
		and since $2(d+m-1)\ge d + m$, we obtain that $\mathbf{a}_{m+1}-\mathbf{a}_m \in I^{d+m}$.
		
		Hence since $A$ is $I$-adically complete, we obtain a limit $\mathbf{b} = \lim_{m\to \infty} \mathbf{a}_m$. Since $\mathbf{a}_m \in I^dA^n$ for all $m$, we have $\mathbf{b}\in I^dA^n$, and taking the limit of $F(\mathbf{a}_m) = \mathbf{a}_{m}-\mathbf{a}_{m+1}$ as $m\to \infty$ we obtain that $F(\mathbf{b}) = 0$, as desired.
	\end{proof}

\end{document}
